% Generated from pinned audit evidence.
\newcommand{\TrainSeconds}{361.6}
\newcommand{\TrainMemory}{0.543}
\newcommand{\FirstIdentity}{1.17}
\newcommand{\FirstShuffled}{1.17}
\newcommand{\FirstZero}{0.00}
\newcommand{\BlockIdentity}{0.586}
\newcommand{\BlockGain}{0.098}
\newcommand{\BlockLow}{0.000}
\newcommand{\BlockHigh}{0.307}
\newcommand{\BlockFone}{0.320}
\newcommand{\RunDecision}{reject}

% Generated from the pinned V54 audit; no fitted presentation values.
\newcommand{\DensitySeconds}{207.8}
\newcommand{\DensityMemory}{0.400}
\newcommand{\DensityNll}{15.219}
\newcommand{\DensityShuffledNll}{15.254}
\newcommand{\DensityAdvantage}{0.0355}
\newcommand{\DensityLow}{-0.1014}
\newcommand{\DensityHigh}{0.1693}
\newcommand{\DensityIdentity}{1.562}
\newcommand{\DensityFone}{0.350}
\newcommand{\DensityCopyFone}{0.99949}
\newcommand{\DensityDecision}{reject}

% Generated from the locked V55 matched audit.
\newcommand{\JointNll}{7.026}
\newcommand{\ControlNll}{7.579}
\newcommand{\JointNllGain}{0.553}
\newcommand{\JointNllLow}{0.450}
\newcommand{\JointNllHigh}{0.665}
\newcommand{\JointTime}{342.4}
\newcommand{\ControlTime}{218.0}
\newcommand{\JointMemory}{0.648}
\newcommand{\ControlMemory}{0.400}
\newcommand{\JointSampleFone}{0.362}
\newcommand{\JointPrefixGain}{0.185}
\newcommand{\JointDecision}{reject}

% Generated from unchanged V56 report.
\newcommand{\BeamAccuracy}{11.621}
\newcommand{\BeamShuffled}{9.863}
\newcommand{\GreedyAccuracy}{8.252}
\newcommand{\BeamGain}{3.369}
\newcommand{\BeamGainLow}{2.351}
\newcommand{\BeamGainHigh}{4.412}
\newcommand{\BeamHistoryGain}{1.758}
\newcommand{\BeamHistoryLow}{0.966}
\newcommand{\BeamHistoryHigh}{2.552}
\newcommand{\BeamRepeat}{25.446}
\newcommand{\TargetRepeat}{0.112}
\newcommand{\BeamSeconds}{269.3}
\newcommand{\BeamMemory}{0.154}
\newcommand{\BeamDecision}{reject\_decoder}

% Generated from the completed V57 audit.
\newcommand{\QaExact}{3.906}
\newcommand{\QaFone}{15.718}
\newcommand{\QaBlankExact}{0.391}
\newcommand{\QaBlankFone}{5.553}
\newcommand{\QaControlExact}{0.000}
\newcommand{\QaControlFone}{6.582}
\newcommand{\QaQuestionGain}{10.165}
\newcommand{\QaQuestionLow}{7.043}
\newcommand{\QaQuestionHigh}{13.371}
\newcommand{\QaPairExact}{0.254}
\newcommand{\QaMotor}{100.000}
\newcommand{\QaSeconds}{278.572}
\newcommand{\QaMemory}{1.053}
\newcommand{\QaControlSeconds}{269.126}
\newcommand{\QaControlMemory}{1.049}
\newcommand{\QaDecision}{reject\_reading\_model}

% Generated from the fixed autonomous verification receipt.
\newcommand{\CodecLanguageFOne}{32.671}
\newcommand{\CodecLanguageBlankFOne}{19.508}
\newcommand{\CodecLanguageQuestionGain}{13.163}
\newcommand{\CodecLanguageCILow}{10.933}
\newcommand{\CodecLanguageCIHigh}{15.730}

\newcommand{\PaperTitle}{Learning conditional language\\in continuous image space}
\newcommand{\PaperSummary}{Written language includes forms outside standard character encodings. We investigate continuous interfaces that read and emit glyph images, with emitted pixels supplying subsequent history. Pretrained foundations provide inherited knowledge; tokenizers, lexical heads and font rendering are absent from the deployed answer path. Two compact fits use 13,785 complete Chinese sequences, 2.576 million target glyph positions and 3.930/4.889 GiB peak allocated GPU memory, but give zero complete exact answers out of 32 at five checkpoints. A wider writer removes a measured rank restriction without solving language. An earlier actor retaining 7.522B visual and contextual parameters gives two fixed-layout and six proportional exact answers; external native decoding gives ten from matched first states. A direct-context motor reaches seven proportional exact answers but fails retention. Supervision through all sampling intervals improves first-position training writing but fails the later-position gate. A physically qualified deterministic writer reproduces all 64 sampled training glyphs exactly, yet gives six proportional exact answers, no both-correct pair and no credited complete book continuation. Its language pilot also fails. Cold replays verify fixed-case independent pixel-feedback execution. These controls distinguish inherited answer access and physical writing from autonomous language transfer. Usable open-ended writing, historical-form transfer and a matched-quality efficiency advantage remain unresolved.}
